\providecommand{\rank}{\operatorname{r}}
\providecommand{\RR}{\mathbb R}
\providecommand{\CC}{\mathbb C}
\providecommand{\KK}{\mathbb K}

% ---------- Exercise 85 ----------
\setcounter{enumi}{84}
\item {\bfseries\boldmath
Consider the vector subspace $F=\langle A_1,A_2,A_3,A_4\rangle$ of $M_2(\RR)$, where
}

\[
A_1=\begin{bmatrix}1&0\\0&0\end{bmatrix},\quad A_2=\begin{bmatrix}1&-1\\1&0\end{bmatrix},\quad A_3=\begin{bmatrix}1&1\\-1&0\end{bmatrix},\quad A_4=\begin{bmatrix}1&1\\1&0\end{bmatrix}.
\]
Determine

\begin{enumerate}[label=\alph*)]

\item a basis of $F$.

\begin{answer}
Let
\[
 A_1=\begin{bmatrix}1&0\\0&0\end{bmatrix},\quad
 A_2=\begin{bmatrix}1&-1\\1&0\end{bmatrix},\quad
 A_3=\begin{bmatrix}1&1\\-1&0\end{bmatrix},\quad
 A_4=\begin{bmatrix}1&1\\1&0\end{bmatrix}.
\]

Since $A_3=2A_1-A_2$, the space is generated by $A_1,A_2,A_4$.
Moreover,
\[
 \alpha A_1+\beta A_2+\gamma A_4
 =\begin{bmatrix}\alpha+\beta+\gamma&-\beta+\gamma\\\beta+\gamma&0\end{bmatrix}.
\]
If this is zero, its off-diagonal entries give $\beta=\gamma=0$, and then
$\alpha=0$. Thus $\mathcal B=(A_1,A_2,A_4)$ is a basis of $F$.
\end{answer}

\item the coordinates of $A=\begin{bmatrix}2&-2\\4&0\end{bmatrix}$ relative to the basis of $F$ found in (a).

\begin{answer}
For $A=\begin{bmatrix}2&-2\\4&0\end{bmatrix}$, compare entries:
\[
 -\beta+\gamma=-2,\quad\beta+\gamma=4,\quad\alpha+\beta+\gamma=2.
\]
Hence $\beta=3$, $\gamma=1$, $\alpha=-2$, so
\[
 A=-2A_1+3A_2+A_4,\qquad [A]_{\mathcal B}=\begin{bmatrix}-2\\3\\1\end{bmatrix}.
\]
\end{answer}
\end{enumerate}
